\appendix
\section{Xu's dense-core lemmas for arbitrary colourings}\label{app:xu}
In this appendix, we prove Lemmas~\ref{lem:critical} and~\ref{lem:remaining}.
Xu~\cite{Xu} proves them for proper colourings, and his proofs go through for arbitrary globally $m$-bounded colourings except at a few steps, where he uses that edges of the same colour do not share a vertex.
We follow his notation and indicate only these steps, giving a replacement argument for each; every other step applies verbatim.
In all of his sampling steps, each large colour class, or each part of it, is sampled by choosing a fixed number of its edges uniformly at random.
The small colour classes are pooled and sampled independently, and different classes are sampled independently.
Two facts about sampling without replacement replace properness.
First, if $s$ elements of a set of size $a$ are chosen uniformly at random, then two given elements are both chosen with probability $s(s-1)/(a(a-1))\le(s/a)^2$.
Second, by Hoeffding's comparison~\cite[Theorem~4]{Hoe}, for any weights on the elements, the total weight of the chosen elements has moment generating function at most that of the total weight of $s$ independent uniform choices.
In particular, the Chernoff bounds used by Xu hold for sums over the chosen edges with arbitrary nonnegative weights, for example for the number of chosen edges at a vertex, even when several of them have the same colour.

\subsection{A spanning tree of a dense core and one more edge}
Lemma~\ref{lem:critical} is deduced from the following lemma, in which $G[C]$ is close to a complete graph and exactly $m$ edges lie outside $C$.

\begin{lemma}[{Xu~\cite[Lemma~5.1]{Xu} for arbitrary colourings}]\label{lem:xu51}
For every $\alpha\in(0,1/2)$ and $\eps>0$ there exist $\delta>0$ and $n_0$ such that the following holds.
Let $G$ have a globally $m$-bounded edge-colouring, and let $C\subseteq V(G)$ have size $n\ge n_0$, where $m\ge\alpha n$.
Suppose that $H=G[C]$ has $(n-1)m$ edges, that exactly $m$ edges of $G$ are not inside $C$, and that $|d_H(v)-2m|\le\delta n$ and $|\,|N_H(u)\cap N_H(v)|-4m^2/n|\le\delta n$ for all distinct $u,v\in C$.
Then $G$ contains $(1-\eps)m$ edge-disjoint rainbow forests, each consisting of a spanning tree of $H$ and one edge not inside $C$.
\end{lemma}
\begin{proof}
We use the notation of~\cite[Section~5]{Xu}: $s\ge\beta n$ is the current number of forests still to be found, $L=(\log n)^{20}$, $p=L/s$ and $D=2n/(\log n)^8$. A colour is large if it has at least $s/(\log n)^2$ edges. The edges of a large colour inside $C$ and those not inside $C$ form two separate parts, and from a part of size $a$ we choose $\lceil pa\rceil$ edges uniformly at random.
Properness is used twice in the sampling step: to show that the support $J$ satisfies $\Delta(J[C])\le D$, and in the bound~\cite[(5.3)]{Xu} on $\sum_w\Prb(uw,vw\in T)$, which is needed for tracking codegrees.

Say that a part inside $C$ of a large colour is \emph{tiny} if $pa<1$. Exactly one edge of a tiny part is chosen. Every other part chooses each of its edges with probability $\lceil pa\rceil/a\le2p$.
A large colour with a tiny part inside $C$ has at least $s/(\log n)^2-s/L\ge s/(2(\log n)^2)$ edges not inside $C$, and only $s$ edges are not inside $C$. So there are at most $2(\log n)^2$ tiny parts, and they have at most $2(\log n)^2s/L=2s/(\log n)^{18}$ edges in total.

For the maximum degree, fix $v\in C$. Tiny parts give at most $2(\log n)^2$ edges of $J$ at $v$. The other chosen edges at $v$ outside the pool kept in full each appear with probability at most $2p$, so their expected number is at most $2p\cdot2.1s=4.2L$. By the Chernoff bound discussed above, there are at most $20L$ of them at every vertex with probability $1-n^{-100}$. The pool kept in full has at most $s/(\log n)^8$ edges. Hence $\Delta(J[C])\le D$.

For~\cite[(5.3)]{Xu}, fix distinct $u,v\in C$ and let $K$ be the pool kept in full. If $uw$ or $vw$ lies in $K$ or in a tiny part, we use $\Prb(uw,vw\in T)\le1/s+n^{-50}$. There are at most $2|K|+4s/(\log n)^{18}\le3s/(\log n)^8$ such $w$, contributing at most $4(\log n)^{-8}$ in total. For every other $w$, both edges are chosen at random. If they lie in different parts or pools, they are chosen independently, each with probability at most $2p$. If they lie in the same part, which can only happen if they have the same colour, then they are both chosen with probability at most $(\lceil pa\rceil/a)^2\le4p^2$, and the same holds in the sampled pool, which is sampled independently. In all cases $\Prb(uw,vw\in J)\le4p^2$, and so $\Prb(uw,vw\in T)\le8p^2$ after conditioning on the likely event in Xu's proof. Summing over at most $n$ vertices $w$ gives at most $8np^2\le8\beta^{-2}(\log n)^{40}/n$. This proves~\cite[(5.3)]{Xu} with $\theta_n=C_\beta((\log n)^{-8}+(\log n)^{40}/n)$, and the rest of Xu's proof applies.
\end{proof}

\begin{proof}[Proof of Lemma~\ref{lem:critical}]
Xu's proof~\cite[Lemma~5.2]{Xu} uses dependent rounding~\cite{GKPS} to choose target numbers $X_\gamma$ and $Y_\gamma$ of edges of each colour $\gamma$ inside and not inside $C$. It then keeps $X_\gamma$ uniformly random edges of colour $\gamma$ inside $C$ and $Y_\gamma$ not inside $C$, independently for different colours.
The resulting graph has $(2m-1)m'$ edges inside $C$ and $m'$ edges outside it, where $m'=m-2W$ and $W\le\eta m$ is the number of edges missing from $G[C]$, and it is globally $m'$-bounded.
None of this uses properness.
Properness is used only to show that every vertex of $C$ loses at most $8W+m^{2/3}$ edges inside $C$, after which Lemma~\ref{lem:xu51} applies as in~\cite{Xu}.
We prove this bound for arbitrary colourings.

Fix $v\in C$. For a colour $\gamma$, let $A_\gamma$ be its number of edges inside $C$, let $d_\gamma\le A_\gamma$ be its number of edges inside $C$ at $v$, and let $D_\gamma$ be the number of these $d_\gamma$ edges which are deleted. Let $D=\sum_\gamma D_\gamma$.
Given the rounding, the $D_\gamma$ are independent, and $D_\gamma$ counts the edges at $v$ among $A_\gamma-X_\gamma$ uniformly random edges of colour $\gamma$ inside $C$.
For $\lambda>0$ let $\varkappa=e^\lambda-1$ and $t_\gamma=1-X_\gamma/A_\gamma$. By Hoeffding's comparison,
\[
\E\bigl[e^{\lambda D_\gamma}\bigm|X_\gamma\bigr]\le(1+\varkappa d_\gamma/A_\gamma)^{A_\gamma-X_\gamma}\le\exp(\varkappa d_\gamma t_\gamma)=:\phi_\gamma(X_\gamma).
\]
In the rounding, $X_\gamma=\lfloor x_\gamma\rfloor+I_\gamma$, where $x_\gamma\ge\rho A_\gamma$ with $\rho=m'/m$, and the indicator variables $I_\gamma$ all lie on edges at a single vertex of the auxiliary bipartite graph. For such variables, dependent rounding gives $\E\prod_{\gamma\in S}(1-I_\gamma)\le\prod_{\gamma\in S}\E(1-I_\gamma)$ for every set $S$ of colours~\cite{GKPS} (see also~\cite[Section~1]{CVZ}).
Since $\phi_\gamma$ is decreasing, $\phi_\gamma=c_\gamma+c'_\gamma(1-I_\gamma)$ for some $c_\gamma,c'_\gamma\ge0$, and expanding the product gives $\E e^{\lambda D}\le\E\prod_\gamma\phi_\gamma\le\prod_\gamma\E\phi_\gamma$.
The variable $d_\gamma t_\gamma$ takes two values, which differ by $d_\gamma/A_\gamma\le1$, and its mean is $d_\gamma(1-x_\gamma/A_\gamma)\le(1-\rho)d_\gamma$. By Hoeffding's lemma, $\E\phi_\gamma\le\exp(\varkappa(1-\rho)d_\gamma+\varkappa^2d_\gamma/8)$.
As $\sum_\gamma d_\gamma\le2m$ and $1-\rho=2W/m$,
\[
\Prb(D\ge8W+m^{2/3})\le\exp\bigl(4\varkappa W+\varkappa^2m/4-\lambda(8W+m^{2/3})\bigr).
\]
With $\lambda=m^{-1/3}$, so that $\varkappa\le\lambda+\lambda^2$, the exponent is at most $-m^{1/3}/2$ for large $m$. A union bound over the $2m$ vertices of $C$ completes the proof.
\end{proof}

\subsection{The remaining case}
\begin{proof}[Proof of Lemma~\ref{lem:remaining}]
In the sampling step of~\cite[Lemma~6.1]{Xu}, each colour with at least $m^b$ edges, where $b=1-\kappa/32$, chooses $\lceil pa_\gamma\rceil$ of its $a_\gamma$ edges uniformly at random, where $p=m^{-1+\kappa/8}$. The small colours form a pool $K$ kept in full if it has at most $m^{1-\kappa/16}$ edges, and otherwise sampled independently with probability $p$.
The weights, the graphic constraints, the support bounds and the scalar tracking use only the edge probabilities and the Chernoff bounds discussed above.
Properness is used only in the matrix tracking. There, Xu needs the bound~\cite[(6.11)]{Xu} that two edges of different colours not in $K$ are both chosen with probability at most $8p^2$, and notes that two edges with a common vertex have different colours.
We show that this bound holds for all pairs of edges not in $K$.
If the two edges are chosen independently, this is as in~\cite{Xu}.
If they have the same colour $\gamma$ and $\gamma$ is large, then they are both chosen with probability at most $(\lceil pa_\gamma\rceil/a_\gamma)^2\le4p^2$, as $pa_\gamma\ge1$.
If they are in the sampled pool, they are chosen independently.
After conditioning on the likely event in Xu's proof, the bound $8p^2$ follows in each case, and the rest of the proof applies verbatim.
\end{proof}
